% Source: https://artofproblemsolving.com/community/c3h2182711_hard_problems_plz_help
% Stable identity: aops-topic-2182711-q02
For two points $\left(p,\dfrac{1}{p}\right)$ and $\left(q,\dfrac{1}{q}\right)$ of the graph with $p\ne q$, the slope of the line through them is
\[\dfrac{\dfrac{1}{q}-\dfrac{1}{p}}{q-p}=\dfrac{\dfrac{p-q}{pq}}{q-p}=-\dfrac{1}{pq},\]
which is never $0$ and never undefined; so two such lines are perpendicular exactly when their slopes multiply to $-1$. Let $C=\left(t,\dfrac{1}{t}\right)$ with $t\ne0,1,3$.
\begin{itemize}
\item Right angle at $A$: $\left(-\dfrac{1}{3}\right)\left(-\dfrac{1}{t}\right)=-1$, so $t=-\dfrac{1}{3}$.
\item Right angle at $B$: $\left(-\dfrac{1}{3}\right)\left(-\dfrac{1}{3t}\right)=-1$, so $t=-\dfrac{1}{9}$.
\item Right angle at $C$: $\left(-\dfrac{1}{t}\right)\left(-\dfrac{1}{3t}\right)=\dfrac{1}{3t^2}>0$, which is never $-1$; impossible.
\end{itemize}

\begin{asy}
usepackage("amsmath");
size(6cm,9cm);
defaultpen(fontsize(10pt));
import graph;
real f(real x){return 1/x;}
draw(graph(f,0.24,4.2),gray);
draw(graph(f,-1.4,-0.105),gray);
draw((-1.6,0)--(4.4,0),Arrow(5)); draw((0,-9.6)--(0,4.4),Arrow(5));
pair A=(1,1), B=(3,1/3), C1=(-1/3,-3), C2=(-1/9,-9);
draw(A--B--C1--cycle,blue);
draw(A--B--C2--cycle,red+dashed);
real s=0.35;
pair u=unit(B-A), v=unit(C1-A); draw(A+s*u--A+s*(u+v)--A+s*v,blue);
pair u2=unit(A-B), v2=unit(C2-B); draw(B+s*u2--B+s*(u2+v2)--B+s*v2,red);
dot(A); dot(B); dot(C1,blue); dot(C2,red);
label("$A(1,1)$",(1.45,1.65),fontsize(9pt)); label("$B\left(3,\dfrac{1}{3}\right)$",(3.35,1.0),fontsize(9pt));
label("$C_1\left(-\dfrac{1}{3},-3\right)$",C1,W,blue+fontsize(9pt));
label("$C_2\left(-\dfrac{1}{9},-9\right)$",C2,E,red+fontsize(9pt));
label("$xy=1$",(0.35,f(0.35)),E,gray+fontsize(9pt));
label("\mbox{}",(min(currentpicture,user=true)+max(currentpicture,user=true))/2,black); // keeps the text after the figure black when LaTeX typesets the labels inline
\end{asy}

So the possible points are $C_1=\left(-\dfrac{1}{3},-3\right)$ and $C_2=\left(-\dfrac{1}{9},-9\right)$, and the sum of their $x$-coordinates is
\[-\dfrac{1}{3}-\dfrac{1}{9}=\boxed{-\dfrac{4}{9}}.\]
